\section*{Hoofdstuk \ref{ch:g10:soil-formation} --- Bodem, een hulpbron die langzaam ontstaat}

\begin{solution}{exo:g10:soil-formation:1}
\omterm{def:g10:soil-formation:chemical-weathering}{Chemische verwering} is de omvorming van de mineralen van een gesteente door
reacties met water en de gassen die het meevoert (koolstofdioxide, zuurstof):
sommige mineralen lossen op, andere worden nieuwe mineralen die aan het oppervlak
stabiel zijn. \omterm{def:g10:soil-formation:hydrolysis}{Hydrolyse} is de \omterm{def:g10:soil-formation:chemical-weathering}{chemische verwering} van een silicaat door water en
zijn waterstofionen: de metaalionen gaan in oplossing en de rest van het mineraal
wordt omgebouwd tot een \omterm{def:g10:soil-formation:clay}{kleimineraal}.
\end{solution}

\begin{solution}{exo:g10:soil-formation:2}
\omterm{def:g10:soil-formation:weathering-profile}{Saproliet} is gesteente dat ter plaatse is ontbonden: zacht en aardachtig, behoudt
het de korrels, banden en structuur van zijn \omterm{def:g10:soil-formation:weathering-profile}{moedergesteente} en gaat het naar
beneden toe, zonder onderbreking, over in het verse gesteente. Een bodem die in
rivierzand is gevormd, rust op materiaal dat van elders is aangevoerd: zijn
korrels zijn afgerond en gesorteerd, en hij toont niets van de structuur van het
gesteente eronder.
\end{solution}

\begin{solution}{exo:g10:soil-formation:3}
De veldspaten worden klei (kaoliniet) en geven daarbij kalium, natrium of
calcium en silica af aan het water; de zwarte mica wordt klei en ijzeroxiden;
het kwarts blijft onveranderd achter, als zand.
\end{solution}

\begin{solution}{exo:g10:soil-formation:4}
Het ijzer van de donkere mineralen, \ce{Fe^2+}, wordt door de zuurstof die het
water meevoert geoxideerd tot onoplosbare ijzer(III)-oxiden: geelbruin goethiet
en rood hematiet, die de korrels bedekken en het hele materiaal kleuren.
\end{solution}

\begin{solution}{exo:g10:soil-formation:5}
\qty{0.017}{mm} per jaar en \qty{1.537}{mm} per jaar: $1.537/0.017 \approx 90$
keer groter.
\end{solution}

\begin{solution}{exo:g10:soil-formation:6}
Calcium: 1 en 1. Koolstof: $1 + 1 = 2$ en 2. Zuurstof: $3 + 2 + 1 = 6$ en
$2 \times 3 = 6$. Waterstof: 2 en 2. Lading: 0, en $+2 - 2 = 0$.
\end{solution}

\begin{solution}{exo:g10:soil-formation:7}
Het oplossen van calciet levert alleen opgeloste ionen op, \ce{Ca^2+} en
\ce{HCO3-}: alles wordt afgevoerd en een zuivere kalksteen laat geen vast residu
achter. De \omterm{def:g10:soil-formation:hydrolysis}{hydrolyse} van veldspaat levert een vaste stof op, kaoliniet, die
achterblijft; en het kwarts van het graniet reageert nauwelijks. Graniet laat
daarom klei en zand achter.
\end{solution}

\begin{solution}{exo:g10:soil-formation:8}
De kleischilfers en de humus dragen negatieve ladingen die de positieve
kaliumionen \ce{K+} op hun oppervlakken vasthouden, zodat het regenwater dat
erdoorheen stroomt ze niet meevoert. Een wortel geeft waterstofionen af, die de
plaats van het kalium op de schilfers innemen; het kalium dat zo in het
bodemwater vrijkomt, wordt door de wortel opgenomen.
\end{solution}

\begin{solution}{exo:g10:soil-formation:9}
\omterm{def:g10:soil-formation:chemical-weathering}{Chemische verwering} heeft water nodig, dat in natte klimaten overvloedig is en in
droge schaars (en in koude bevroren); en haar reacties verlopen sneller wanneer
het warm is. Alleen bij warmte en overvloedig water gaan de reacties, over lange
tijd, ver genoeg om bijna alles wat oplosbaar is te verwijderen en ijzer- en
aluminiumoxiden achter te laten.
\end{solution}

\begin{solution}{exo:g10:soil-formation:10}
$10/0.017 \approx 590$ jaar voor \qty{1}{cm}; $300/0.017 \approx 17\,600$
jaar voor \qty{30}{cm}; $1000/0.017 \approx 59\,000$ jaar voor \qty{1}{m}.
\end{solution}

\begin{solution}{exo:g10:soil-formation:11}
Volume: $10\,000 \times 0.0015 = \qty{15}{m^3}$; massa: $15 \times 1200 =
\qty{18000}{kg}$, \num{18}~ton per hectare per jaar.
\end{solution}

\begin{solution}{exo:g10:soil-formation:12}
$46\,000/0.017 \approx 2.7$ miljoen jaar. Zo'n dikke \omterm{def:g10:soil-formation:weathering-profile}{saproliet} vraagt een
landschap dat miljoenen jaren stabiel is gebleven, zonder dat zijn verweerde
mantel werd weggeruimd.
\end{solution}

\begin{solution}{exo:g10:soil-formation:13}
Goethiet: \ce{4 Fe^2+ + O2 + 6 H2O -> 4 FeO(OH) + 8 H+}: ijzer 4 en 4,
zuurstof $2 + 6 = 8$ en $4 \times 2 = 8$, waterstof 12 en $4 + 8 = 12$, lading
$+8$ en $+8$. Hematiet: \ce{4 Fe^2+ + O2 + 4 H2O -> 2 Fe2O3 + 8 H+}: ijzer 4
en 4, zuurstof $2 + 4 = 6$ en 6, waterstof 8 en 8, lading $+8$ en $+8$.
\end{solution}

\begin{solution}{exo:g10:soil-formation:14}
Nettoverlies $0.082 - 0.017 = \qty{0.065}{mm}$ per jaar; $300/0.065 \approx 4600$
jaar. Bij conventioneel ploegen $300/1.52 \approx 200$ jaar: meer dan twintig
keer korter.
\end{solution}

\begin{solution}{exo:g10:soil-formation:15}
De ton bevat \qty{300}{kg} orthoklaas. Elke kilogram geeft \qty{464}{g}
kaoliniet: $300 \times 0.464 \approx \qty{139}{kg}$. Kalium:
$300 \times 0.1405 \approx \qty{42}{kg}$.
\end{solution}

\begin{solution}{pb:g10:soil-formation:1}
\textbf{1.} $4 \times 0.3048 \approx \qty{1.22}{m}$.

\textbf{2.} $0.5 \times 30.48 \approx \qty{15}{cm}$.

\textbf{3.} Humus, het donkere organische materiaal dat afbrekers uit dode
planten maken.

\textbf{4.} IJzeroxiden, goethiet en hematiet, ontstaan door de oxidatie van het
ijzer(II) van de mineralen van het moedermateriaal.

\textbf{5.} Nee. Daaronder zou het profiel overgaan in verweerd gesteente, de
\omterm{def:g10:soil-formation:weathering-profile}{saproliet} (of ander verweerd materiaal), en dan in het verse \omterm{def:g10:soil-formation:weathering-profile}{moedergesteente}.

\textbf{6.} \ce{2 KAlSi3O8 + 2 H+ + 9 H2O -> Al2Si2O5(OH)4 + 2 K+ + 4 H4SiO4}.

\textbf{7.} $1000/278.33 \approx \qty{3.59}{mol}$.

\textbf{8.} $3.59/2 \approx \qty{1.80}{mol}$ kaoliniet, $1.80 \times 258.16
\approx \qty{464}{g}$.

\textbf{9.} De kaliumionen en het opgeloste silica: het water voert ze naar de
rivieren, en de rivieren voeren ze naar zee.

\textbf{10.} Het reageert nauwelijks: de kwartskorrels blijven als zand achter.

\textbf{11.} $1219/0.017 \approx 72\,000$ jaar.

\textbf{12.} $1.537 - 0.017 = \qty{1.52}{mm}$ per jaar.

\textbf{13.} $1219/1.52 \approx 800$ jaar.

\textbf{14.} $72\,000/800 = 90$: het profiel gaat negentig keer sneller verloren
dan het is ontstaan.

\textbf{15.} Nee: $0.013 < 0.017$, de erosie is iets trager dan de productie en
de bodem houdt stand.

\textbf{16.} $T = S/(E - P)$.

\textbf{17.} $300/1.52 \approx 200$ jaar.

\textbf{18.} $300/(0.082 - 0.017) = 300/0.065 \approx 4600$ jaar.

\textbf{19.} Een bodem die in een paar eeuwen verloren gaat, heeft tienduizenden
jaren nodig om opnieuw te ontstaan: op de schaal van mensenlevens en
samenlevingen wordt wat verloren gaat niet vervangen.

\textbf{20.} $1000/1.52 \approx 660$ jaar: \textbf{bij conventioneel ploegen gaat
een meter bodem in ongeveer 660 jaar verloren, een paar eeuwen}.
\end{solution}
